\chapter{Conception et réalisation -- Volet Intelligence Artificielle}

\section*{Introduction}
Ce troisième chapitre constitue le cœur technique du rapport. Il présente l'architecture globale du système, les différents diagrammes UML qui en formalisent la conception, la modélisation de la base de données, ainsi que le détail complet du volet Intelligence Artificielle~: constitution des jeux de données, prétraitement, entraînement et évaluation des modèles de Machine Learning, puis leur intégration dans l'application.

\section{Architecture globale}

GREENLY repose sur une architecture en trois niveaux (\textit{3-tier}), clairement séparée entre présentation, logique métier/IA et persistance~:

\begin{itemize}[label=\textbullet,font=\normalsize]
    \item une \textbf{application front-end} \textit{Single Page Application} (SPA) développée en React et TypeScript, responsable de l'interface utilisateur~;
    \item un \textbf{micro-service de Machine Learning} indépendant, développé en Python avec le framework FastAPI, exposant une API REST dédiée aux prédictions, aux prévisions et aux métriques de modèles~;
    \item une \textbf{base de données} PostgreSQL hébergée sur la plateforme Supabase, assurant à la fois la persistance des données utilisateurs, l'authentification et le contrôle d'accès.
\end{itemize}

Cette séparation permet de faire évoluer indépendamment la partie prédictive (ré-entraînement, ajout d'algorithmes) et la partie applicative, et facilite le déploiement de chaque composant dans son propre conteneur Docker (voir le diagramme de déploiement, figure~\ref{fig:deploiement}). La figure~\ref{fig:composants} présente le diagramme de composants de l'architecture globale.

\begin{figure}[H]
\centering
\resizebox{0.95\textwidth}{!}{
\begin{tikzpicture}[>=Stealth, font=\small]

\tikzset{
  comp/.style={draw, rounded corners=2pt, minimum width=4.6cm, minimum height=1.3cm, align=center, line width=0.7pt},
  stereo/.style={font=\scriptsize\itshape}
}

% Client
\node[comp, fill=isiBlue!8] (client) at (0,6) {\stereo{«device»}\\\textbf{Navigateur Web (Client)}};

% Frontend
\node[comp, fill=isiBlue!15, minimum height=2.6cm, text width=5cm] (frontend) at (0,3) {
  \stereo{«component»}\\
  \textbf{Frontend SPA}\\
  \footnotesize React 18 + TypeScript + Vite\\
  Pages, Services, Hooks, Contexts
};

% ML Service
\node[comp, fill=isiBlue!25, minimum height=2.6cm, text width=5cm] (mlservice) at (-4.5,-0.5) {
  \stereo{«component»}\\
  \textbf{Micro-service ML}\\
  \footnotesize FastAPI + scikit-learn\\
  Routers, ModelRegistry, Pipelines
};

% Supabase
\node[comp, fill=isiBlue!25, minimum height=2.6cm, text width=5cm] (supabase) at (4.5,-0.5) {
  \stereo{«component»}\\
  \textbf{Supabase}\\
  \footnotesize PostgreSQL + Auth + RLS\\
  Table \texttt{analyses}
};

\draw[<->] (client) -- (frontend) node[midway, right, font=\scriptsize] {HTTP(S)};
\draw[<->] (frontend) -- (mlservice) node[midway, above, sloped, font=\scriptsize] {REST/JSON};
\draw[<->] (frontend) -- (supabase) node[midway, above, sloped, font=\scriptsize] {SDK JS};

\node[font=\scriptsize, align=left, text width=4.6cm] at (-4.5,-2.6) {\texttt{/api/predict/\{category\}}\\\texttt{/api/forecast}\\\texttt{/api/models/metrics}};
\node[font=\scriptsize, align=left, text width=4.6cm] at (4.5,-2.6) {CRUD table \texttt{analyses}\\Authentification (JWT)\\Row Level Security};

\end{tikzpicture}
}
\caption{Diagramme de composants -- architecture globale de GREENLY}
\label{fig:composants}
\end{figure}

\section{Diagrammes UML}

\subsection{Diagramme de classes}

Le diagramme de classes de la figure~\ref{fig:classes} modélise le domaine du micro-service de Machine Learning~: les schémas d'entrée par catégorie (validés par Pydantic), le registre de modèles \texttt{ModelRegistry} qui orchestre les prédictions, la réponse structurée \texttt{PredictionResponse}, ainsi que les classes dédiées au module de prévision.

\begin{figure}[H]
\centering
\resizebox{\textwidth}{!}{
\begin{tikzpicture}[>=Stealth, font=\scriptsize]

\tikzset{
  cls/.style={draw, rectangle split, rectangle split parts=3, text width=5cm, align=left, line width=0.6pt},
}

\node[cls] (input) at (0,0) {
  \centering\textbf{CalculatorInput} \scriptsize(6 variantes)
  \nodepart{second}
  + vehicleType, fuelType, distance, ... \newline
  + electricityKwh, heatingType, ... \newline
  + showerMinutes, laundryLoads, ... \newline
  + meatMeals, dairyServings, ... \newline
  + plasticKg, paperKg, recycledPercent, ... \newline
  + deviceType, usageHours, powerWatts, ...
  \nodepart{third}
  \textit{(validation Pydantic~: bornes min/max)}
};

\node[cls, text width=5.6cm] (registry) at (9,0) {
  \centering\textbf{ModelRegistry}
  \nodepart{second}
  - pipelines~: dict[str, Pipeline] \newline
  - metrics~: dict \newline
  - distributions~: dict
  \nodepart{third}
  + predict(category, input)~: dict \newline
  + percentile\_rank(category, value)~: float \newline
  + is\_ready(category)~: bool
};

\node[cls, text width=5.6cm] (response) at (9,-5.3) {
  \centering\textbf{PredictionResponse}
  \nodepart{second}
  + category, prediction, target, unit \newline
  + greenScore, confidence, modelUsed \newline
  + percentileRank, featureContributions
  \nodepart{third}
  \textit{(structure de retour JSON)}
};

\node[cls, text width=5cm] (pipeline) at (0,-5.3) {
  \centering\textbf{Pipeline (scikit-learn)}
  \nodepart{second}
  - preprocessor~: ColumnTransformer \newline
  - model~: RandomForest / GBM / \newline~~~XGBoost / LightGBM / CatBoost
  \nodepart{third}
  + fit(X, y) \newline
  + predict(X)~: array
};

\node[cls, text width=5cm] (forecast) at (0,-10.6) {
  \centering\textbf{ForecastEngine}
  \nodepart{second}
  (module fonctionnel)
  \nodepart{third}
  + forecast(points, periodsAhead)~: dict
};

\node[cls, text width=5.6cm] (forecastresp) at (9,-10.6) {
  \centering\textbf{ForecastResponse}
  \nodepart{second}
  + points~: List[ForecastResponsePoint] \newline
  + trend, trendPercent \newline
  + modelConfidence, method
  \nodepart{third}
  \textit{(structure de retour JSON)}
};

\draw[->] (input) -- node[above, font=\scriptsize] {1..1} node[below, font=\scriptsize]{utilise} (registry);
\draw[->] (registry) -- node[left, font=\scriptsize] {compose} (pipeline);
\draw[->] (registry) -- node[right, font=\scriptsize] {crée} (response);
\draw[->] (forecast) -- node[below, font=\scriptsize] {crée} (forecastresp);
\draw[->] (registry.south) .. controls (4.5,-8) .. (forecast.north) node[midway, above, font=\scriptsize]{partage \texttt{distributions.json}};

\end{tikzpicture}
}
\caption{Diagramme de classes -- domaine du micro-service Machine Learning}
\label{fig:classes}
\end{figure}

\subsection{Diagramme de séquence}

La figure~\ref{fig:sequence} illustre le scénario nominal \textit{« un utilisateur effectue un calcul d'empreinte carbone »}, depuis la saisie du formulaire jusqu'à l'enregistrement du résultat dans l'historique.

\begin{figure}[H]
\centering
\resizebox{\textwidth}{!}{
\begin{tikzpicture}[>=Stealth, font=\scriptsize]

\def\lifelines{Utilisateur/0, Frontend (React)/3.2, API FastAPI/7, ModelRegistry/10.6, Supabase/14}

\foreach \name/\x in \lifelines {
    \node[draw, fill=isiBlue!15, minimum width=2.6cm, minimum height=0.8cm, align=center, line width=0.6pt] (head-\x) at (\x,0) {\name};
    \draw[dashed] (\x,0) -- (\x,-13.5);
}

% Messages (from-x, to-x, y, label, style)
\draw[->] (0,-1) -- (3.2,-1) node[midway, above, font=\scriptsize]{1. Saisit le formulaire};
\draw[->] (3.2,-1.9) -- (3.4,-1.9) node[midway, above, font=\scriptsize, xshift=1.4cm]{2. Valide localement (bornes)};
\draw[->] (3.2,-2.8) -- (7,-2.8) node[midway, above, font=\scriptsize]{3. POST /api/predict/\{category\}};
\draw[->] (7,-3.7) -- (7.2,-3.7) node[midway, above, font=\scriptsize, xshift=1.9cm]{4. Valide le schéma (Pydantic)};
\draw[->] (7,-4.6) -- (10.6,-4.6) node[midway, above, font=\scriptsize]{5. predict(category, data)};
\draw[->] (10.6,-5.5) -- (10.8,-5.5) node[midway, above, font=\scriptsize, xshift=1.9cm]{6. pipeline.predict(row)};
\draw[->, dashed] (10.6,-6.4) -- (7,-6.4) node[midway, above, font=\scriptsize]{7. prediction, greenScore, ...};
\draw[->, dashed] (7,-7.3) -- (3.2,-7.3) node[midway, above, font=\scriptsize]{8. 200 OK (JSON)};
\draw[->] (3.2,-8.2) -- (3.4,-8.2) node[midway, above, font=\scriptsize, xshift=1.4cm]{9. Affiche le résultat};
\draw[->] (3.2,-9.4) -- (14,-9.4) node[midway, above, font=\scriptsize]{10. insert analyses (inputs, outputs, co2e)};
\draw[->] (14,-10.3) -- (14.2,-10.3) node[midway, above, font=\scriptsize, xshift=2cm]{11. Vérifie RLS (auth.uid() = user\_id)};
\draw[->, dashed] (14,-11.2) -- (3.2,-11.2) node[midway, above, font=\scriptsize]{12. Confirmation};
\draw[->] (3.2,-12.4) -- (0.3,-12.4) node[midway, above, font=\scriptsize]{13. Affiche la confirmation};

\end{tikzpicture}
}
\caption{Diagramme de séquence -- réalisation d'un calcul d'empreinte}
\label{fig:sequence}
\end{figure}

\subsection{Diagramme d'activité}

La figure~\ref{fig:activite} détaille le déroulement du pipeline d'entraînement des modèles de Machine Learning (script \texttt{train.py}), exécuté hors ligne pour chacune des six catégories.

\begin{figure}[H]
\centering
\begin{tikzpicture}[>=Stealth, font=\scriptsize, node distance=0.55cm and 1cm]

\tikzset{
  act/.style={draw, rounded corners=3pt, fill=isiBlue!10, minimum width=6.4cm, minimum height=0.9cm, align=center, line width=0.6pt},
  startstop/.style={circle, draw, fill=black, minimum size=0.28cm},
  endstop/.style={circle, draw, minimum size=0.4cm}
}

\node[startstop] (start) at (0,0) {};
\node[act] (a1) at (0,-1.1) {Charger le jeu de données brut réel (CSV public)};
\node[act] (a2) at (0,-2.3) {Nettoyer et filtrer les données};
\node[act] (a3) at (0,-3.5) {Construire la variable cible (formule physique documentée + bruit réaliste)};
\node[act] (a4) at (0,-4.9) {Encoder les variables catégorielles (\textit{One-Hot Encoding})};
\node[act] (a5) at (0,-6.1) {Diviser en jeu d'entraînement (80\%) et de test (20\%)};
\node[act, minimum height=1.3cm] (a6) at (0,-7.7) {Entraîner et évaluer les 5 algorithmes\\(Random Forest, Gradient Boosting, XGBoost, LightGBM, CatBoost)};
\node[act] (a7) at (0,-9.3) {Sélectionner le modèle au R\textsuperscript{2} le plus élevé};
\node[act] (a8) at (0,-10.5) {Sauvegarder le pipeline retenu (\texttt{.joblib})};
\node[act] (a9) at (0,-11.7) {Générer \texttt{metrics.json} et \texttt{distributions.json}};
\node[endstop] (end) at (0,-12.7) {};
\node[circle, draw, fill=black, minimum size=0.18cm] at (end) {};

\draw[->] (start) -- (a1);
\draw[->] (a1) -- (a2);
\draw[->] (a2) -- (a3);
\draw[->] (a3) -- (a4);
\draw[->] (a4) -- (a5);
\draw[->] (a5) -- (a6);
\draw[->] (a6) -- (a7);
\draw[->] (a7) -- (a8);
\draw[->] (a8) -- (a9);
\draw[->] (a9) -- (end);

\node[font=\scriptsize, align=left, text width=4.5cm] at (5.2,-7.7) {\textit{Répété pour chacune des 6 catégories~: carbon, energy, electronics, water, food, waste}};
\draw[->, dashed] (3.3,-7.7) -- (5.2,-7.7);

\end{tikzpicture}
\caption{Diagramme d'activité -- pipeline d'entraînement des modèles}
\label{fig:activite}
\end{figure}

\subsection{Diagramme de déploiement}

La figure~\ref{fig:deploiement} présente l'architecture de déploiement telle que définie dans le fichier \texttt{docker-compose.yml}~: deux conteneurs Docker orchestrés (frontend et micro-service ML), communiquant avec la plateforme Supabase hébergée dans le cloud.

\begin{figure}[H]
\centering
\resizebox{0.92\textwidth}{!}{
\begin{tikzpicture}[>=Stealth, font=\scriptsize]

\tikzset{
  node3d/.style={draw, rounded corners=2pt, minimum width=5cm, minimum height=1.6cm, align=center, line width=0.7pt},
  stereo/.style={font=\scriptsize\itshape}
}

\node[node3d, fill=isiBlue!8] (browser) at (0,4) {\stereo{«device»}\\\textbf{Poste client}\\Navigateur Web};

\node[node3d, fill=isiBlue!20, minimum height=2.1cm] (nginx) at (-4.5,0) {
  \stereo{«conteneur Docker»}\\
  \textbf{frontend}\\
  \footnotesize nginx:1.27-alpine\\
  Port 8080 → 80\\
  Artefact~: \texttt{dist/} (build React)
};

\node[node3d, fill=isiBlue!20, minimum height=2.1cm] (mlnode) at (4.5,0) {
  \stereo{«conteneur Docker»}\\
  \textbf{ml-service}\\
  \footnotesize python:3.x + FastAPI + uvicorn\\
  Port 8000\\
  Artefacts~: \texttt{*\_model.joblib}
};

\node[node3d, fill=isiBlue!35, minimum height=1.9cm] (supa) at (0,-4) {
  \stereo{«service Cloud»}\\
  \textbf{Supabase}\\
  \footnotesize PostgreSQL + Auth + RLS
};

\draw[<->] (browser) -- (nginx) node[midway, left, font=\scriptsize]{HTTP :8080};
\draw[<->] (browser) -- (mlnode) node[midway, right, font=\scriptsize]{REST/JSON :8000};
\draw[<->] (nginx) -- (supa) node[midway, below, sloped, font=\scriptsize]{HTTPS (SDK)};
\draw[<->, dashed] (nginx) -- (mlnode) node[midway, above, font=\scriptsize]{depends\_on: healthy};

\end{tikzpicture}
}
\caption{Diagramme de déploiement -- environnement Docker Compose}
\label{fig:deploiement}
\end{figure}

\section{Base de données}

GREENLY utilise une base de données PostgreSQL hébergée sur Supabase. Le schéma repose sur une table principale, \texttt{analyses}, reliée à la table système \texttt{auth.users} gérée nativement par Supabase pour l'authentification. La figure~\ref{fig:erd} présente le diagramme entité-association correspondant.

\begin{figure}[H]
\centering
\begin{tikzpicture}[>=Stealth, font=\scriptsize]

\node[draw, rectangle split, rectangle split parts=2, text width=4.6cm, align=left, line width=0.7pt, fill=isiBlue!15] (users) at (0,0) {
  \centering\textbf{auth.users} \scriptsize(géré par Supabase)
  \nodepart{second}
  \underline{id}~: uuid (PK)\\
  email~: text\\
  created\_at~: timestamptz
};

\node[draw, rectangle split, rectangle split parts=2, text width=6.2cm, align=left, line width=0.7pt, fill=isiBlue!8] (analyses) at (9,0) {
  \centering\textbf{analyses}
  \nodepart{second}
  \underline{id}~: uuid (PK)\\
  user\_id~: uuid (FK → auth.users.id)\\
  category~: text \textit{(check~: 6 valeurs)}\\
  inputs~: jsonb\\
  outputs~: jsonb\\
  ai\_recommendation~: text\\
  green\_score~: integer\\
  co2e~: numeric\\
  notes~: text\\
  created\_at~: timestamptz\\
  updated\_at~: timestamptz
};

\draw[-] (users.east) -- (analyses.west) node[midway, above, font=\scriptsize]{1} node[near end, above, font=\scriptsize]{N};

\end{tikzpicture}
\caption{Diagramme entité-association -- schéma de la base de données}
\label{fig:erd}
\end{figure}

Un utilisateur (\texttt{auth.users}) peut posséder de zéro à plusieurs analyses (\texttt{analyses}), relation modélisée par la clé étrangère \texttt{user\_id}. Trois index sont définis pour optimiser les requêtes les plus fréquentes~: sur \texttt{user\_id} (filtrage par utilisateur), sur \texttt{created\_at} (tri chronologique de l'historique) et sur \texttt{category} (filtrage par type de calculateur). La sécurité d'accès est assurée par quatre politiques \textbf{RLS} (\textit{Row Level Security}) -- une par opération CRUD -- garantissant qu'un utilisateur authentifié ne peut lire, insérer, modifier ou supprimer que les lignes dont il est propriétaire~:

\begin{lstlisting}
CREATE POLICY "select_own_analyses" ON analyses FOR SELECT
  TO authenticated USING (auth.uid() = user_id);
\end{lstlisting}

\section{Architecture Frontend / Backend}

\subsection{Frontend}
L'interface utilisateur est développée en \textbf{React 18} avec \textbf{TypeScript}, à l'aide de l'outil de build \textbf{Vite}. La mise en forme repose sur \textbf{Tailwind CSS} combiné aux composants accessibles \textbf{Radix UI} (via la bibliothèque de composants shadcn/ui). Le routage est assuré par \texttt{react-router-dom}, les formulaires par \texttt{react-hook-form} couplé à des schémas de validation \texttt{zod}, et les visualisations par la bibliothèque \texttt{recharts}. L'organisation du code suit une architecture en couches~:

\begin{itemize}[label=\textbullet,font=\normalsize]
    \item \texttt{pages/} -- les vues de l'application (landing, authentification, calculateurs, tableaux de bord)~;
    \item \texttt{services/} -- la couche d'accès aux données externes (\texttt{mlService.ts} pour l'API ML, \texttt{analysisService.ts} pour Supabase, \texttt{pdfReportService.ts} pour l'export PDF)~;
    \item \texttt{lib/} -- la logique métier pure (\texttt{calculators.ts}, formatage, configuration des calculateurs)~;
    \item \texttt{hooks/} et \texttt{contexts/} -- la gestion d'état partagé (authentification, thème, historique des analyses)~;
    \item \texttt{components/} -- les composants d'interface réutilisables.
\end{itemize}

\subsection{Backend -- micro-service Machine Learning}
Le micro-service est développé en Python avec \textbf{FastAPI}, choisi pour sa validation native des données via \textbf{Pydantic}, ses performances asynchrones et la génération automatique d'une documentation interactive (OpenAPI/Swagger). Il est structuré en trois modules principaux~: \texttt{routers/} (points d'entrée REST), \texttt{registry.py} (chargement et orchestration des modèles), et \texttt{schemas.py} (contrats de données d'entrée/sortie).

\subsection{Backend -- persistance}
La persistance et l'authentification sont déléguées à \textbf{Supabase}, une plateforme \textit{Backend-as-a-Service} bâtie sur PostgreSQL. Ce choix permet de bénéficier d'une authentification robuste, d'une sécurité déclarative (RLS) et d'un SDK JavaScript directement intégrable côté client, sans développer de backend applicatif dédié à la gestion des utilisateurs.

\section{Description des modules}

Le tableau~\ref{tab:modules} résume les principaux modules du micro-service Machine Learning et leur responsabilité.

\begin{table}[H]
\centering
\caption{Description des modules du micro-service ML}
\label{tab:modules}
\begin{tabular}{|p{3.6cm}|p{9.5cm}|}
\hline
\rowcolor{isiBlue!15}
\textbf{Module} & \textbf{Responsabilité} \\
\hline
\texttt{scripts/download\_datasets.py} & Télécharge les jeux de données publics bruts (EPA, UCI, Our World in Data). \\ \hline
\texttt{scripts/prepare\_datasets.py} & Nettoie les données brutes et construit les variables cibles (labels) d'entraînement selon des formules physiques documentées. \\ \hline
\texttt{scripts/train.py} & Entraîne et compare 5 algorithmes par catégorie, sélectionne et sauvegarde le meilleur modèle. \\ \hline
\texttt{app/schemas.py} & Définit les schémas Pydantic de validation des requêtes et réponses de l'API. \\ \hline
\texttt{app/registry.py} & Charge les pipelines entraînés en mémoire et expose la fonction de prédiction. \\ \hline
\texttt{app/forecasting.py} & Calcule la tendance statistique (régression linéaire) sur l'historique d'un utilisateur. \\ \hline
\texttt{app/routers/predict.py} & Expose les routes \texttt{POST /api/predict/\{category\}}. \\ \hline
\texttt{app/routers/analytics.py} & Expose les routes de métriques de modèles et de prévision. \\ \hline
\texttt{app/main.py} & Point d'entrée FastAPI~: assemblage des routers, configuration CORS. \\ \hline
\end{tabular}
\end{table}

\section{Dataset}

Contrairement à une approche à formule fixe, GREENLY fonde ses prédictions sur l'entraînement de modèles supervisés. Or, il n'existe pas de jeu de données public unique reliant directement des habitudes de consommation individuelles à une empreinte carbone mesurée. Une approche \textbf{hybride} a donc été retenue~: pour chaque catégorie, la relation entre les variables d'entrée et la valeur cible est ancrée dans des données publiques réelles et documentées, complétées -- lorsque nécessaire -- par une simulation de distributions d'usage réalistes.

Le tableau~\ref{tab:datasets} présente, pour chaque catégorie, la source réelle utilisée et la nature du jeu de données d'entraînement obtenu.

\begin{table}[H]
\centering
\caption{Sources de données par catégorie}
\label{tab:datasets}
\begin{tabular}{|p{2.1cm}|p{4.6cm}|p{2cm}|p{4.4cm}|}
\hline
\rowcolor{isiBlue!15}
\textbf{Catégorie} & \textbf{Source réelle} & \textbf{Taille} & \textbf{Nature} \\
\hline
Transport & EPA \texttt{fueleconomy.gov} -- \texttt{vehicles.csv} (émissions mesurées réelles) \cite{epaVehicles} & 67~189 lignes & Réelle + distances simulées \\ \hline
Énergie & UCI \textit{Appliances Energy Prediction} \cite{candanedo2017} (mesures de capteurs réelles) & 12~000 lignes & Réelle + mise à l'échelle par foyer \\ \hline
Électronique & Constantes EPA ENERGY STAR / DOE (puissances typiques) \cite{energyStar} & 8~000 lignes & Simulation à partir de constantes réelles \\ \hline
Eau & Constantes EPA WaterSense / USGS (débits typiques) \cite{waterSense} & 8~000 lignes & Simulation à partir de constantes réelles \\ \hline
Alimentation & Our World in Data / Poore \& Nemecek (2018) \cite{poore2018} (émissions réelles par aliment) & 8~000 lignes & Réelle (intensités) + repas simulés \\ \hline
Déchets & Our World in Data -- taux de collecte des déchets \cite{owidWaste} & 8~000 lignes & Réelle (taux de recyclage) + composition simulée \\ \hline
\end{tabular}
\end{table}

Chaque constante physique utilisée (facteur d'émission du réseau électrique, intensité carbone d'un type de chauffage, débit d'une douche économe, etc.) est documentée dans le code source et rattachée à sa source officielle, garantissant la traçabilité scientifique du système -- un point que le chapitre~1 identifiait comme une lacune des outils existants.

\section{Prétraitement}

Le prétraitement des données se déroule en deux temps~: la construction du jeu de données d'entraînement (variable cible), puis sa préparation pour l'algorithme d'apprentissage.

\subsection{Construction de la variable cible}
Pour la catégorie Transport, à titre d'exemple, l'émission mesurée par l'EPA (en grammes de CO2 par mile, valeur réelle) est convertie en kilogrammes par kilomètre, puis combinée à une distance simulée et à un facteur correctif « conditions réelles vs laboratoire » (écart documenté par l'EPA et l'ICCT, de l'ordre de 10 à 25\,\%)~:

\begin{lstlisting}[language=Python]
co2_per_km = (co2TailpipeGpm / 1000) / 1.60934
real_world_factor = np.clip(RNG.normal(1.15, 0.08), 0.9, 1.45)
co2e_kg = co2_per_km * distance_km * real_world_factor
\end{lstlisting}

\subsection{Encodage et découpage}
Les variables catégorielles (par exemple \texttt{vehicleType}, \texttt{fuelType}) sont encodées par \textbf{One-Hot Encoding}, tandis que les variables numériques sont conservées telles quelles (les modèles à base d'arbres de décision ne nécessitant pas de normalisation). Le jeu de données est ensuite scindé en un ensemble d'entraînement (80\,\%) et un ensemble de test (20\,\%), ce dernier n'étant jamais utilisé pendant l'apprentissage afin de garantir une évaluation non biaisée~:

\begin{lstlisting}[language=Python]
preprocessor = ColumnTransformer([
    ("cat", OneHotEncoder(handle_unknown="ignore"), categorical),
    ("num", "passthrough", numeric),
])
X_train, X_test, y_train, y_test = train_test_split(
    X, y, test_size=0.2, random_state=42
)
\end{lstlisting}

\section{Entraînement des modèles de Machine Learning}

Pour chacune des six catégories, cinq algorithmes de régression supervisée sont entraînés et comparés automatiquement~:

\begin{itemize}[label=\textbullet,font=\normalsize]
    \item \textbf{Random Forest} \cite{breiman2001}~: forêt de 300 arbres de décision (profondeur maximale~14), moyennant leurs prédictions individuelles.
    \item \textbf{Gradient Boosting}~: 300 arbres construits séquentiellement, chacun corrigeant l'erreur résiduelle du précédent (taux d'apprentissage~0,05).
    \item \textbf{XGBoost} \cite{chen2016}~: implémentation optimisée du \textit{gradient boosting} (400 arbres, profondeur~6).
    \item \textbf{LightGBM} \cite{ke2017}~: variante de \textit{gradient boosting} à croissance par feuille, optimisée pour la vitesse (400 arbres).
    \item \textbf{CatBoost} \cite{prokhorenkova2018}~: \textit{gradient boosting} avec un traitement natif optimisé des variables catégorielles (400 itérations, profondeur~6).
\end{itemize}

Ces cinq algorithmes reposent tous sur le principe des \textbf{ensembles d'arbres de décision}, particulièrement adaptés aux données tabulaires mixtes (catégorielles et numériques) et capables de modéliser des interactions non linéaires entre variables, contrairement aux formules linéaires fixes utilisées par les outils concurrents (voir chapitre~1). Le meilleur algorithme est sélectionné automatiquement selon le coefficient de détermination R\textsuperscript{2} obtenu sur le jeu de test~:

\begin{lstlisting}[language=Python]
for algo_name, make_model in ALGORITHMS.items():
    pipeline = Pipeline([("prep", preprocessor), ("model", make_model())])
    pipeline.fit(X_train, y_train)
    preds = pipeline.predict(X_test)
    r2 = r2_score(y_test, preds)
    if r2 > best_score:
        best_name, best_score, best_pipeline = algo_name, r2, pipeline
joblib.dump(best_pipeline, MODELS / f"{category}_model.joblib")
\end{lstlisting}

\section{Évaluation}

Trois métriques standard de régression sont utilisées pour évaluer chaque modèle sur le jeu de test (données jamais vues pendant l'entraînement)~:

\begin{itemize}[label=\textbullet,font=\normalsize]
    \item \textbf{MAE} (\textit{Mean Absolute Error})~: erreur absolue moyenne, dans l'unité de la variable cible~;
    \item \textbf{RMSE} (\textit{Root Mean Squared Error})~: pénalise davantage les grosses erreurs~;
    \item \textbf{R\textsuperscript{2}} (coefficient de détermination)~: proportion de la variance de la cible expliquée par le modèle (1 = prédiction parfaite).
\end{itemize}

Le tableau~\ref{tab:comparaison-algos} présente la comparaison complète des 5 algorithmes pour la catégorie Transport, obtenue lors de l'exécution réelle du pipeline d'entraînement.

\begin{table}[H]
\centering
\caption{Comparaison des 5 algorithmes -- catégorie Transport (67~189 trajets, 20\,\% test)}
\label{tab:comparaison-algos}
\begin{tabular}{|l|c|c|c|}
\hline
\rowcolor{isiBlue!15}
\textbf{Algorithme} & \textbf{MAE (kg)} & \textbf{RMSE (kg)} & \textbf{R\textsuperscript{2}} \\
\hline
\textbf{Gradient Boosting (retenu)} & \textbf{1,878} & \textbf{4,218} & \textbf{0,9824} \\ \hline
Random Forest & 1,902 & 4,629 & 0,9788 \\ \hline
LightGBM & 1,980 & 6,561 & 0,9574 \\ \hline
CatBoost & 2,033 & 6,854 & 0,9535 \\ \hline
XGBoost & 2,017 & 7,058 & 0,9507 \\ \hline
\end{tabular}
\end{table}

Le tableau~\ref{tab:resultats-globaux} synthétise le modèle finalement retenu pour chacune des six catégories du système.

\begin{table}[H]
\centering
\caption{Modèle retenu et performance par catégorie (données réelles \texttt{metrics.json})}
\label{tab:resultats-globaux}
\begin{tabular}{|l|l|c|c|c|}
\hline
\rowcolor{isiBlue!15}
\textbf{Catégorie} & \textbf{Algorithme retenu} & \textbf{R\textsuperscript{2}} & \textbf{MAE} & \textbf{RMSE} \\
\hline
Énergie & CatBoost & 0,9993 & 1,808 & 2,550 \\ \hline
Déchets & CatBoost & 0,9993 & 0,052 & 0,104 \\ \hline
Électronique & CatBoost & 0,9875 & 0,036 & 0,061 \\ \hline
Transport & Gradient Boosting & 0,9824 & 1,878 & 4,218 \\ \hline
Eau & CatBoost & 0,7820 & 57,76 & 77,38 \\ \hline
Alimentation & CatBoost & 0,5667 & 11,85 & 15,76 \\ \hline
\end{tabular}
\end{table}

L'analyse de ces résultats appelle une remarque importante~: la précision varie sensiblement selon la catégorie. Les catégories Énergie, Déchets et Électronique, où la cible dépend directement de facteurs physiques appliqués aux variables observées, atteignent une précision quasi parfaite. À l'inverse, les catégories Eau et surtout Alimentation intègrent, lors de la construction du jeu d'entraînement, des facteurs de variation aléatoire (par exemple la composition exacte du régime alimentaire d'un utilisateur) qui ne sont \textbf{pas} observables dans les variables d'entrée fournies au modèle~; ce bruit non observable constitue un plafond théorique de précision, quel que soit l'algorithme utilisé -- et non un défaut de modélisation.

\section{Intégration des modèles IA dans l'application}

Les pipelines entraînés (prétraitement + modèle, sérialisés au format \texttt{.joblib}) sont chargés une seule fois, au démarrage du micro-service, par la classe \texttt{ModelRegistry} (voir diagramme de classes, figure~\ref{fig:classes}), qui les conserve en mémoire pour garantir des temps de réponse de l'ordre de la dizaine de millisecondes. Chaque requête de prédiction déclenche la séquence suivante (détaillée dans le diagramme de séquence, figure~\ref{fig:sequence})~:

\begin{enumerate}
    \item validation du corps de la requête par un schéma Pydantic dédié à la catégorie~;
    \item construction d'une ligne de données au format attendu par le pipeline~;
    \item appel de \texttt{pipeline.predict()}, qui applique successivement le prétraitement (encodage) puis le modèle entraîné~;
    \item calcul du \textit{Green Score} par positionnement percentile de la prédiction dans la distribution des valeurs observées à l'entraînement (\texttt{distributions.json})~;
    \item retour d'une réponse structurée incluant la prédiction, le score, le niveau de confiance (R\textsuperscript{2} du modèle) et l'importance des variables.
\end{enumerate}

Ce découplage entre entraînement (hors ligne, via \texttt{scripts/train.py}) et inférence (en ligne, via l'API FastAPI) permet de ré-entraîner ou d'ajouter de nouveaux algorithmes sans jamais interrompre le service applicatif.

\section*{Conclusion}
Ce chapitre a détaillé la conception et la réalisation du volet Intelligence Artificielle de GREENLY~: de l'architecture logicielle globale jusqu'à l'intégration effective des modèles entraînés dans l'application, en passant par la constitution rigoureuse des jeux de données, leur prétraitement et l'évaluation comparative de cinq algorithmes d'apprentissage supervisé. Les résultats obtenus, avec un coefficient R\textsuperscript{2} atteignant 0,999 pour certaines catégories, démontrent la viabilité de l'approche. Le chapitre suivant présente le second volet du projet~: le module de Business Intelligence et de prévision, qui exploite ces modèles pour offrir à l'utilisateur un suivi analytique de son impact environnemental.
